% Generated draft: requires mathematical and source review.
\documentclass[11pt]{article}
\usepackage[margin=25mm]{geometry}
\usepackage{fontspec}
\setmainfont{Latin Modern Roman}
\usepackage{amsmath,amssymb,mathtools,unicode-math}
\setmathfont{Latin Modern Math}
\usepackage{enumitem,needspace,xurl,hyperref}
\setlength{\parindent}{0pt}
\setlength{\parskip}{5pt}
\setlength{\emergencystretch}{3em}
\newcommand{\N}{\mathbb{N}}
\newcommand{\Z}{\mathbb{Z}}
\newcommand{\Q}{\mathbb{Q}}
\newcommand{\R}{\mathbb{R}}
\newcommand{\C}{\mathbb{C}}
\newcommand{\F}{\mathbb{F}}
\newcommand{\A}{\mathbb{A}}
\newcommand{\PP}{\mathbb{P}}
\newcommand{\E}{\mathbb{E}}
\newcommand{\eps}{\varepsilon}
\newcommand{\dd}{\,\mathrm{d}}
\begin{document}
\section*{UnsolvedMath ID: 30003797}
\textbf{Source problem code:} catalogue entry
\subsection{Definitions and mathematical statement}
A field \(K\) is finitely generated if it is finitely generated as a field extension of its prime field. An absolutely almost simple algebraic \(K\)-group \(G\) is a nontrivial connected semisimple algebraic group whose base change to an algebraic closure of \(K\) has no proper nontrivial connected normal algebraic subgroup.

The supplied statement requires \(K\) to have ``good characteristic,'' but does not define that requirement or specify its dependence on \(G\). It also does not specify the ambient family of places. Write \(V\) for that intended family, let \(S\subseteq V\) be finite, and write \(K_v\) for the completion at \(v\), whenever the intended notion of place supplies such a completion. These are placeholders for missing source specifications, not additional choices resolving them.

For a field extension \(L/K\), let \(H^1(L,G)\) denote the nonabelian Galois cohomology pointed set of \(G_L\), equivalently the set of isomorphism classes of \(G_L\)-torsors, with distinguished point the trivial torsor. Extension of scalars gives the localization map
\[
\operatorname{loc}_{V\setminus S}\colon H^1(K,G)\longrightarrow\prod_{v\in V\setminus S}H^1(K_v,G).
\]
Define its distinguished fiber by
\[
\mathcal{X}(G)_S
:=\operatorname{loc}_{V\setminus S}^{-1}\bigl((1_v)_{v\in V\setminus S}\bigr)
=\bigl\{\xi\in H^1(K,G):\xi|_{K_v}=1_v\text{ for every }v\in V\setminus S\bigr\},
\]
where \(1_v\) is the class of the trivial \(G_{K_v}\)-torsor. This is a pointed set; no group structure is assumed.

The supplied finiteness question asks whether, for every finitely generated field \(K\) and every absolutely almost simple \(K\)-group \(G\) satisfying the intended good-characteristic hypothesis, and every finite \(S\) in the intended family \(V\), one has
\[
\#\mathcal{X}(G)_S<\infty.
\]
The hypothesis on characteristic and the family \(V\) must be recovered before this becomes a fully specified conjecture. In particular, this draft does not assert the claim for arbitrary families of places. The original statement additionally asks whether ``Question 2'' has a positive answer, with the displayed finiteness question as a particular consequence. The content of Question 2 is absent from the supplied record, so that broader target cannot be formulated here.
\subsection{Short English statement}
The question asks whether there are only finitely many isomorphism classes of \(G\)-torsors over \(K\) that become trivial over every specified completion outside a finite exceptional set \(S\). The intended places, the meaning of good characteristic, and the broader referenced Question 2 remain unspecified.
\subsection{Sources}
\begin{itemize}
\item Supplied local UnsolvedMath registry record, permanent ID \(30003797\), source problem code \texttt{catalogue entry}, title: \emph{Finiteness of Tate--Shafarevich Sets for Algebraic Groups}. Collection: Oberwolfach Reports Open Problems. No individual proposer is supplied; the proposed year is recorded as 2018.
\item Supplied report citation: \emph{Quadratic Forms and Related Structures over Fields (2019)}, Oberwolfach Reports, DOI 10.4171/owr/2018/21, pp.~1235--1291. \url{https://doi.org/10.4171/owr/2018/21}. The background attributes the problem to workshop open-problem discussions documented in Oberwolfach Reports.
\item Supplied literature source, labeled ``spinor groups with good reduction'': \url{https://arxiv.org/abs/1707.08062}.
\item Supplied literature source, labeled ``good-reduction survey'': \url{https://arxiv.org/abs/2005.05484}.
\item Supplied literature source, labeled ``workshop report'': \url{https://ems.press/content/serial-article-files/46742}.
\item Literature assessment dated exactly \texttt{2026-08-21}: the supplied curation labels the problem \texttt{partially\_solved} and states that the Tate--Shafarevich finiteness conjecture follows in known type-A, spinor, and related cases, while the general finitely generated good-characteristic conjecture remains unresolved. This is supplied dated curation derived from the OWR material, not an independently verified assessment; its time-sensitive status and cited claims require independent verification.
\item The registry labels the statement \texttt{corrected\_verified} and reports that it was checked against the official Oberwolfach report, preserving its hypotheses, notation, and open target. This is a supplied verification claim; no source consultation or mathematical verification is asserted by this draft.
\end{itemize}
\end{document}
